\section{Robustness}\label{sec:robustness}

\subsection{Restricting to the Balanced Subsample}\label{subsec:balanced}

Our baseline results pool balanced and unbalanced individuals by absorbing the latter into a single pooled cell and interacting this group dummy with urban status.
Appendix~\ref{app:pooling} shows that this pooling is consistent under our identifying assumptions plus two additional conditions: \textit{(i)} comparative advantage in the ``unbalanced'' group does not vary systematically across sector, once we've conditioned on $x_{i t}$, and \textit{(ii)} the partial effect of covariates on log consumption is the same for the balanced and unbalanced groups.
Below, we nonetheless verify empirically that the baseline estimates do not depend on the unbalanced subsample by re-estimating the restricted GRC on the balanced subsample.
% \footnote{The close correspondence between pooled and balanced-only trajectory estimates provides an empirical check on the joint null that the orthogonality and common-covariate-effects conditions of Appendix~\ref{app:pooling} both hold.}

\GRCtable{IDN}{consumption}{urban}{bal}
\GRCtable{CHN}{consumption}{urban}{bal}
\GRCtable{TZA}{consumption}{urban}{bal}

Tables \ref{tab:GRC_IDN_consumption_urban_bal}, \ref{tab:GRC_CHN_consumption_urban_bal}, and \ref{tab:GRC_TZA_consumption_urban_bal}  show the results for Indonesia, China, and Tanzania, respectively.
The results for China and Tanzania, shown in Tables \ref{tab:GRC_CHN_consumption_urban_bal} and \ref{tab:GRC_TZA_consumption_urban_bal}, are nearly identical to those obtained in the full sample.
The results for Indonesia, shown in Table \ref{tab:GRC_IDN_consumption_urban_bal}, are also similar to those estimated with the unbalanced panel, albeit slightly attenuated. 
In particular, the estimated returns to urban migration for non-migrants fall to 7 log points once we include all controls and a time trend. 
This estimate is very similar to the one that we obtain using fixed effects in the unbalanced panel (column 7 in Table \ref{tab:OLS_cons_urb_unb}). 
However, a more appropriate comparison would be the fixed-effects model estimated on a balanced panel (see column 7 in Table \ref{tab:OLS_cons_urb_bal}), which are much smaller at 3 log points. 

The length of the time period over which the Indonesian data were collected leads to a drastic reduction in the available sample size, which is why our preferred estimates are those with the full sample. 
Comparing the summary statistics across the full and balanced samples, with the latter reported in Table \ref{tab:summary_stats_IDN_bal}, we can see that the age profile of respondents differs quite substantially across the two samples.
In the full sample, rural respondents are on average around one year younger than their urban counterparts. 
In contrast, this difference is reversed in the full sample, with urban respondents being nearly two years older. 
This may reflect different mortality patterns, which would also skew attrition patterns, reinforcing our preference for the full sample. 


%==============================================================================
%  Robustness subsection: cluster-pooling robustness check.
%  Included from main.tex under \section{Robustness}.
%==============================================================================

\subsection{Robustness to cluster pooling}\label{subsec:verdier-robustness}

The restricted GRC of Section~\ref{subsec:restricted-grc-model} identifies $\phi$ from variation across switcher trajectories, with workers in the same trajectory pooled regardless of their cluster of origin.
\citet{verdierAverageTreatmentEffects2020a} observes that this pooling step can pick up bias if cluster assignment is correlated with trajectory: workers in different switcher trajectories who live in different clusters bring different cluster-level consumption levels into the trajectory means, and the differences distort the LCA slope.
Verdier's solution is a worker-level estimator that recovers $\phi$ without aggregating across clusters within trajectory.
We use his cluster-residualization idea to examine whether this is a concern in our data.

We re-estimate $\phi$ from the moment system of Equation~\eqref{eq:restricted-GRC} after replacing the switcher-treatment instruments $D_{it}\cdot\mathbbm 1\{\underline d_i = \underline d\}$ with their cluster residuals.
For each switcher trajectory $\underline d \in \mathcal D_S$ the residualized instrument is
\begin{equation}\label{eq:residualized-instrument}
\tilde D_{it}^{\,\underline d} \;\equiv\; D_{it}\,\mathbbm 1\{\underline d_i=\underline d\} \;-\; \overline{D_{\cdot\cdot}\,\mathbbm 1\{\underline d_{\cdot}=\underline d\}}_{v_i,\,\underline d},
\end{equation}
where the second term is the within-cluster mean of $D_{it}\cdot\mathbbm 1\{\underline d_i=\underline d\}$ taken among trajectory-$\underline d$ workers in cluster $v_i$.
Identification of $\phi$ comes from how much the within-cluster mean of $D$ varies across switcher trajectories, with between-cluster variation absorbed.
Standard errors are clustered at the level of $v_i$, which is the location where we observe person $i$ in the first wave of the data.
In Indonesia and China, we measure location by province, and in Tanzania we use region.
We denote this estimator $\hat\phi^{\mathrm{cluster}}$ and reserve $\hat\phi$ for the baseline of Section~\ref{sec:grc-returns}.

Table~\ref{tab:verdier-robust} reports $\hat\phi^{\mathrm{cluster}}$ alongside the baseline $\hat\phi$ for our preferred consumption specification, which includes all our main control variables and time fixed effects.
Tables X Y Z report the results for all specifications. 


% Generated by cluster_comparison_table (0_programs.do) --- do not hand-edit.
\begin{table}[htb!] \centering \begin{threeparttable}
\caption{Baseline versus cluster-residualized GRC, consumption (unbalanced panel)}
\label{tab:verdier-robust}
\begin{tabular}{l cc cc} \toprule
                       & \multicolumn{2}{c}{$\phi$}
                       & \multicolumn{2}{c}{$\Delta_{\text{never}}$} \\
\cmidrule(lr){2-3} \cmidrule(lr){4-5}
                       & Baseline & Cluster & Baseline & Cluster \\
\midrule
Indonesia & -0.525*** & -0.329** & 0.071*** & 0.057** \\
 & (0.102) & (0.147) & (0.018) & (0.025) \\
\addlinespace
China & -0.205 & -0.157 & 0.098*** & 0.096*** \\
 & (0.133) & (0.231) & (0.021) & (0.035) \\
\addlinespace
China (rural-first) & -0.039 & -0.117 & 0.106*** & 0.104*** \\
 & (0.153) & (0.286) & (0.023) & (0.037) \\
\addlinespace
China (urban-first) & -0.973*** & -0.961*** & 0.006 & 0.031 \\
 & (0.169) & (0.155) & (0.054) & (0.096) \\
\addlinespace
Tanzania & -0.719*** & -0.690*** & 0.270*** & 0.259*** \\
 & (0.124) & (0.087) & (0.033) & (0.033) \\
\bottomrule
\end{tabular}
\begin{tablenotes}[flushleft] \footnotesize
\item{Each row reports two quantities from the restricted GRC on the unbalanced consumption sample with full controls and period fixed effects (column~(5) of the per-country GRC tables): the LCA slope $\phi$ and the extrapolated never-migrant return $\Delta_{\text{never}}$. The baseline columns reproduce the pooled estimates; the cluster columns re-estimate after replacing the switcher-treatment instruments with their within-cluster residuals, using onestep GMM, with standard errors clustered at the level of the cluster of first observation (first-wave province in Indonesia and China, region in Tanzania). Stars: $^{*}p<0.10$; $^{**}p<0.05$; $^{***}p<0.01$.}
\end{tablenotes} \end{threeparttable} \end{table}


The cluster-residualized slope stays negative and close to the baseline in all three countries.
In Tanzania the two are nearly identical, $-0.719$ against $-0.690$, both significant at the 1\% level.
In China both are small and statistically insignificant, $-0.205$ against $-0.155$, so residualizing the instruments changes little.
In Indonesia the slope attenuates from $-0.525$ to $-0.334$ but stays negative and significant at the 5\% level, a gap of about one standard error.
We read this as evidence that pooling switcher trajectories across clusters does not drive the negative $\phi$: the sign survives in every country, and the significance survives in two of three, once between-cluster variation is removed from the identifying instruments.

We read this comparison as a probe of the pooling step, not as a second identification argument.
If workers sort into clusters independently of their switcher trajectory, then residualizing the instruments within cluster leaves $\phi$ unchanged, and $\hat\phi^{\mathrm{rob}}$ tracks the baseline $\hat\phi$.
A material gap between the two would instead signal that cluster composition varies systematically across switcher trajectories and pulls on the pooled slope.
The check holds the identifying restrictions (A1)--(A5) of Section~\ref{subsec:empirical-model} fixed and varies only how the switcher-treatment instruments enter, so it isolates the cluster-pooling margin rather than re-testing (A1)--(A5).
We do not implement the joint-GMM extension of \citet{verdierAverageTreatmentEffects2020a} that simultaneously estimates trajectory-specific average treatment effects for stayers, since our object of interest is $\phi$.


% \subsection{Additional covariates}\label{subsec:xs}

%  \subsection{Homogeneity of Returns Over Time}\label{subsec:shorter_panels}

%  \subsection{Weak identification test}\label{subsec:mu_graph}

% \clearpage

% \subsection{Experience}\label{subsec:experience}
\FloatBarrier
